\subsection{癫痫与性功能}
\label{subsec:epilepsy_sex}

\noindent\textit{【编者注】}癫痫是常见的慢性神经系统疾病。它对性功能的影响涉及疾病本身、抗癫痫药物、心理与社会因素多个层面，却常被患者与医生共同忽视。

\vspace{0.5em}
\noindent\textbf{癫痫如何影响性功能}：

\begin{itemize}
	\item \textbf{疾病本身}：颞叶癫痫（尤其累及边缘系统）可影响性欲、性唤起与性行为；部分患者出现性欲减退，少数表现为性欲亢进或性行为异常（如"露阴自动症"，属器质性疾病表现，并非道德问题）
	\item \textbf{抗癫痫药物的影响}：酶诱导类抗癫痫药（如卡马西平、苯妥英、苯巴比妥）可加速性激素代谢、升高性激素结合球蛋白，导致\textbf{游离睾酮下降}，引起性欲减退、ED、射精障碍与精子质量下降；丙戊酸与多囊卵巢、月经紊乱相关
	\item \textbf{神经内分泌因素}：癫痫发作与脑区放电可干扰下丘脑--垂体--性腺轴的调节
	\item \textbf{心理与社会因素}：病耻感、对发作的恐惧、被歧视的经历、伴侣与家庭的过度保护，都会显著抑制性欲与亲密行为——其影响往往超过疾病本身
	\item \textbf{共病}：抑郁、焦虑在癫痫患者中常见，本身即抑制性欲
\end{itemize}

\vspace{0.5em}
\noindent\textbf{癫痫发作与性生活的安全问题}：

\begin{itemize}
	\item \textbf{癫痫控制优先}：绝大多数患者经规范治疗可正常生活与性生活；关键是坚持服药、定期复诊，\textbf{切勿因担心性功能副作用而自行停药}——发作复发的代价远大于性功能影响
	\item \textbf{体位与安全}：癫痫控制不佳者，可在床边适当留出空间、避免过度疲劳与饮酒诱因、选择双方有安全感的体位；若发作表现为意识丧失，建议避免在水中、高处或危险环境进行性活动
	\item \textbf{发作时的处理}：伴侣应了解患者发作时的正确应对（保护头部、侧卧位、勿强行按压或塞入异物），而不是惊慌失措
	\item \textbf{生育与遗传咨询}：有生育计划者，应接受神经科与生殖医学的联合咨询，评估用药致畸风险与调整方案（如补充叶酸、调整药物）
\end{itemize}

\vspace{0.5em}
\noindent\textbf{评估与处理}：

\begin{enumerate}
	\item 主动向神经科医生说明性功能方面的困扰——很多患者羞于开口，而医生也常不主动询问
	\item 评估可能的药物因素：检查性激素水平（总睾酮、游离睾酮、SHBG、LH/FSH、催乳素），必要时由神经科与内分泌科协商调整药物
	\item 针对性处理：ED 可考虑 PDE5 抑制剂（需注意与部分抗癫痫药的相互作用与低血压风险）；性欲低下需排查抑郁与药物因素
	\item 心理与伴侣支持：减少病耻感、鼓励开放沟通，必要时寻求性治疗或心理治疗
\end{enumerate}

\begin{tcolorbox}[colback=pink!5!white,colframe=red!50!black,title=贴心小叮咛]
癫痫患者完全可以拥有正常的性生活与生育能力。真正伤害性健康的，往往不是癫痫本身，而是\textbf{"不敢问、不好意思说"的沉默}与\textbf{"怕副作用而自行停药"的误解}。把性功能问题带进诊室，与医生一起在"控制发作"和"保护性功能"之间找到平衡，才是正确的路径。
\end{tcolorbox}
